\section{Related Work}
\label{sec:relatedwork}

\subsection{Introduction to Related Works}
This section surveys published work across five areas relevant to the
design and evaluation of the challenge suite: CTF pedagogy and its
integration into graded coursework, progressive difficulty and
multi-stage exploit design, web application vulnerability taxonomies
(particularly OWASP coverage), containerised challenge distribution,
and the emerging use of LLMs as automated CTF solvers. For each area,
the review identifies key findings, evaluates methodological
limitations, and highlights gaps that this project addresses. The
overarching aim is to establish the academic and practitioner
foundations for the design decisions described in the Methodology.


\subsection{CTF Pedagogy and Experiential Learning}
Capture-the-Flag (CTF) competitions have become widely adopted as a
useful pedagogical tool in cybersecurity education. Leune and Petrilli \cite{KeesPetrelli} demonstrate that CTF participation significantly enhances technical proficiency, student engagement, and self-confidence. By providing a hands-on environment for skill refinement, where exercises can effectively bridge the gap between theoretical knowledge and real-world cybersecurity practice. This is aligned with Dewey's principle that practical and theoretical instruction must proceed together. However, the study draws primarily on self-reported student attitudes rather than controlled experimental designs, limiting the strength of their causal claims.

Meinsma et al. \cite{Meinsma2022} provide a measurement framework linking Jeopardy-style challenges to specific knowledge, skill, and attitude objectives using the CTFd platform \cite{CTFd}. Their results challenge the assumption that participation automatically builds competence, finding no clear evidence of measurable learning growth or a consistent correlation between performance and enthusiasm. They concluded that logging data alone cannot reveal how or why learning occurs. For this project, their work underscores the necessity of considering explicit learning objectives within each challenge and capturing process-based evidence of skill development rather than relying solely on flag completion.

The tension between these positions highlights a gap in empirical evidence for learning gain within graded, web-specific coursework. Vykopal et al. \cite{Vykopal} address this by identifying four risks when integrating CTFs into formal assessments: improper scoring, insufficient scaffolding, plagiarism, and limited analytics. They notably found that the pressure of a graded module can shift student focus from conceptual learning to "flag hunting," where the desire to win eclipses the pedagogical intent. This is directly relevant to the current project, as it underscores the need to prioritize instructional design over competitive dynamics to ensure that solving a challenge requires true understanding rather than just guesswork.

The literature establishes a consensus on the pedagogical potential of CTFs, though empirical validation for graded, web-specific coursework remains limited. Effective challenge design frequently employs a multi-level progression and practitioner-driven maxims—such as realism, fairness, and ensuring that learning is the direct path to the flag such as the ones proposed by \cite{Balaji2025}.

The integration of these principles is best understood through Csíkszentmihályi’s Flow Theory \cite{Csikszentmihalyi}, which suggests that optimal engagement occurs when challenge difficulty precisely matches participant skill. This project adopts a tiered structure (from basic cookie tampering to complex SSRF pivoting) to maintain this "flow," mitigating the frustration of over-taxing novices or the boredom of under-challenging experts.

This project addresses these gaps and aligns with the literature by mapping challenges directly to OWASP \cite{OWASP2025} categories and utilizing a dual-evaluation approach—human participant testing and automated LLM analysis—to rigorously assess difficulty and pedagogical alignment.



\subsection{Designing CTF Challenges: Progressive Difficulty and Multi-Stage Exploits}
Balaji et al. \cite{Balaji2025} describe a structured learning path in which each flag unlocks the next level, progressing from basic to advanced topics. This sequential gating mechanism ensures that players demonstrate
competence at each stage before proceeding, which aligns with the
pedagogical principle of scaffolding. However, their platform spans
cryptography, reverse engineering, forensics, and web exploitation
simultaneously, meaning the difficulty gradient operates across domains
rather than within a single domain. The present project adopts both that approach along with a
different approach: the difficulty progression is achieved through increasing the complexity of the exploit chain within the same domain. This within-domain
progression, from a single-step cookie manipulation (CTF1) through to a two-stage JWT forgery (CTF2) allows for more linear scaffolding.

Fuzyll's practitioner maxims \cite{Fuzyll2015} argue that challenges should reward
understanding over guesswork, and that learning should be the path to
the flag rather than a coincidental byproduct. This principle motivates
a deliberate design decision in this project: each challenge embeds
multiple discovery chains leading to the same exploit, so that players
who approach the problem from different angles can still reach the
intended learning outcome. The trade-off is that multiple discovery
paths increase the risk of unintended solutions, which must be
mitigated through post-design vulnerability auditing, a process
documented for each challenge in the project's workflow artefacts A 'safer example' would be in allowing a few edge cases of SQL injection rather than only a singular payload working.


\subsection{Web Application Vulnerabilities and OWASP Coverage}
The selection of vulnerability classes for the challenge suite is
grounded in the OWASP Top 10 \cite{OWASP2025}, the most widely adopted taxonomy of
web application security risks. OWASP serves as a practical framework
for both industry practitioners and educators, and several authors have
used it to structure CTF curricula. Sadat et al. \cite{sadat2024} survey common web
application vulnerabilities alongside countermeasure strategies,
confirming that classic flaws such as SQLi, XSS, and broken authentication remain the most prevalent
targets for defensive training. Nawrocki et al. \cite{nawrocki} provide a
complementary OWASP-mapped analysis, including an experimental test
application and developer checklist. Both papers are useful for
justifying the inclusion of foundational vulnerability classes.

The challenge suite in this project is designed to cover a broad range
of OWASP Top 10 (2021-2025) categories, including Broken Access
Control, Cryptographic Failures, Injection, Identification and Authentication Failures, and Server-Side
Request Forgery. Within the injection category alone, the suite
encompasses SQL injection with filter bypass, DOM-based XSS, and
server-side template injection, each implemented in a distinct
technology stack. This breadth of coverage is deliberately wider than
existing educational CTF platforms, which Sasiwanit et al. \cite{Sasiwanit} observe
typically remain limited to classical vulnerabilities such as brute
force, SQL injection, and XSS, omitting modern exploit classes such as
SSTI, deserialisation, and prototype pollution.

Two vulnerability classes warrant particular discussion because they are under-represented in CTF education literature and form key components of the challenge suite.

Server-Side Template Injection (SSTI) falls under OWASP's Injection category
and occurs when user input is unsafely embedded into a template engine's rendering pipeline. Unlike SQL injection, which targets the database layer, SSTI targets the application's view layer and can escalate directly to remote code execution (RCE) by traversing the template
engine's object hierarchy \cite{PortSwiggerSSTI}. Check
Point Research \cite{CheckPointSSTI2024} reports that as of 2024, an average of one in sixteen organisations faced SSTI attacks weekly, with cloud-based organisations experiencing 30\% more frequent attacks than on-premises environments. Despite this prevalence, SSTI is rarely included in educational CTF platforms. An SSTI challenge in this project such as CTF5 addresses this gap by implementing a Jinja2/Flask environment where players progress from information disclosure through basic template injection to WAF-bypassed RCE, modelling the escalation path that practitioners encounter in real-world engagements.

Server-Side Request Forgery (SSRF) was in the 2021 OWASP Top 10 \cite{OWASP2025}, driven by community consensus regarding its high severity in cloud-native architectures despite a relatively low incidence rate of 2.72\%. The 2019 Capital One breach, analysed by Khan et al. \cite{khan2022}, exemplifies the catastrophic potential of this vulnerability. In that instance, an attacker exploited an SSRF flaw to query the AWS Instance Metadata Service (IMDSv1), obtaining role credentials that exposed over 100 million records. Khan et al. highlighted a critical design flaw in IMDSv1: the service implicitly trusted any local process request, a vulnerability AWS subsequently mitigated with the session-based token authentication of IMDSv2. Literature such as Khan et al. show the importance of the inclusion of SSRF vulnerabilities.


\subsection{Containerisation and Secure Challenge Distribution}
Modern CTF deployments increasingly look into the use Docker containers rather than full virtual machines, offering advantages in scalability, resource
efficiency, and isolation. Raj et al. \cite{raj2016} introduced a CTF
infrastructure using Docker containers, showing that containerisation
significantly reduces the resource requirements for organisers and
automates many setup tasks. Because containers share the host OS kernel,
running many instances is far more efficient than spawning multiple VMs,
and a single machine can host challenges for many participants
simultaneously. Sasiwanit et al. \cite{Sasiwanit} validate this approach in their
CTFd-integrated platform, using Docker for challenge isolation and VPN access control, and demonstrate through Apache JMeter stress testing that the architecture remains responsive under load.
 
However, Neither Raj et al. nor Sasiwanit et al. discuss the
problem of having flag data within the docker images, as their platforms assume a centralised hosting model where students connect to a remote service rather than running containers
locally.

\subsection{LLM-Based Evaluation of CTF Challenges}
Ji et al. \cite{ji2025} reported that LLMs possess substantial technical knowledge about
vulnerability classes but struggle to apply that knowledge to specific
scenarios and to adapt their strategies based on feedback from the CTF
environment. With appropriate scaffolding, including retrieval-augmented
generation and interactive environmental augmentation, their CTFAgent
framework achieved over 80\% performance improvement on standard CTF
datasets. This result has a direct implication for this project: the
difficulty an LLM experiences with a challenge depends heavily on
whether it is given raw prompts or an agentic framework with tool
access, and any LLM evaluation methodology must control for this
variable.


\subsection{Summary of Findings}
The literature establishes broad consensus that CTF exercises are
pedagogically valuable for cybersecurity education, yet rigorous
empirical evidence for measurable learning gain remains limited,
particularly for web-specific challenges embedded within graded
university coursework rather than stand-alone competitions. Existing
educational CTF platforms predominantly cover classical vulnerability
classes such as SQL injection and XSS, with modern exploit categories
including SSTI and SSRF rarely represented despite their growing
real-world prevalence and severity. The problem of securely
distributing containerised challenges while protecting embedded
sensitive data is a practical concern not addressed by the current
literature on Docker-based CTF infrastructure. Finally, LLM-based
evaluation of CTF challenges is a nascent but rapidly developing field;
early benchmarks demonstrate that LLMs struggle with multi-stage
exploitation requiring adaptive reasoning, though performance is highly
dependent on the scaffolding provided. This project addresses these gaps
by designing a challenge suite with broad OWASP coverage including
under-taught vulnerability classes, packaging challenges as
distributable Docker containers, and employing a dual evaluation
methodology combining human participant testing with systematic
LLM-based difficulty assessment.

